% Pipeline choices of scripts/subhalo_multiband/subhalo_multiband_real_data_slam.py (real data) and
% scripts/subhalo_multiband/subhalo_multiband_mock_data_slam.py (mock data). Every number below is the
% value in the code as committed; update both together.

This section lists every choice made in the lens-modelling pipeline, stage by stage. It describes the
real-data pipeline; Section~\ref{sc:pipeline_mock} lists what differs for the mock data. The pipeline
is a Source, Light and Mass (SLaM) chain fitted with \texttt{PyAutoLens}~\citep{Nightingale2021},
followed by the subhalo search of \citet{Nightingale2024}. Every stage is sampled with the nested
sampler \texttt{Nautilus}~\citep{nautilus}.

\subsection{\label{sc:pipeline_general}General set-up}

\paragraph{Joint multi-band fitting.} Each waveband is one factor of a factor graph and all bands are
fitted simultaneously by one non-linear search per stage. A prior object placed in several bands'
models is a single parameter of the joint fit. The same code fits one band or several, so the
single-band and joint-band arms of the experiment differ only in their data.
\begin{itemize}
    \item \textbf{Shared by all bands:} the lens mass (\texttt{Isothermal}, later a
    \texttt{PowerLaw}), the external shear and the dark matter subhalo, because lensing is achromatic.
    \item \textbf{Per band:} the lens light (its own multi-Gaussian expansion, MGE, with its own
    centre and ellipticity, because the stellar light need not have the same centroid or shape at
    every wavelength), the source (an MGE, later a pixelized source with its own mesh and
    regularization), the adapt images, the over-sampling, and the band's alignment to the
    reference band.
\end{itemize}

\paragraph{Reference band and alignment.} The first band given is the astrometric reference. Its
alignment is fixed (zero offset, zero rotation). Every other band has three free alignment
parameters in every stage in which the alignment is free: a $(y, x)$ grid offset, each with a uniform
prior on $\pm 0.5''$, and a grid rotation angle with a uniform prior on $\pm 10^\circ$. The rotation
is about the offset point. The background sky level is fixed at zero in every band. Alignment can be
fixed for all bands with \texttt{--fixed\_alignment}.

\paragraph{Prior chaining.} When a later stage takes a component from an earlier one as a
\emph{model} (free), the installed \texttt{PyAutoFit} reuses the earlier stage's \emph{original}
prior unchanged; the prior is not narrowed by the earlier posterior. A component taken as an
\emph{instance} is fixed at the earlier stage's maximum-likelihood value. The linear light profiles
of a fixed lens-light MGE keep their geometry (centre, ellipticity), but their amplitudes are still
solved linearly in every likelihood evaluation.

\paragraph{Data preparation.} Each band is loaded with its own pixel scale and given a circular mask
of radius $3.5''$ (the \texttt{mask\_radius} of the band's \texttt{info.json}), centred on the
band's \texttt{mask\_centre} (default $(0'', 0'')$). Light profiles are over-sampled in radial bins
about the mask centre: sub-grid size 4 within $0.3''$, 2 between $0.3''$ and $0.6''$, and 2 beyond.
Pixelized sources use the sparse (w-tilde) inversion operator to fit in GPU memory; it gives the same
likelihood as the dense formalism. For real data, every band is plotted and inspected for
contaminating galaxies, stars and artefacts before fitting, and its mask is set from that inspection.

\paragraph{Position penalty.} The multiple-image positions of the lensed source are supplied in the
reference band's frame. The penalty is attached to the reference band's factor only. It constrains
the shared mass model, so one copy suffices, and a copy per band would penalize single- and
multi-band fits differently. Where the threshold is computed from an earlier result, it is twice the
largest source-plane separation of the positions traced through that result's maximum-likelihood
model, clipped to $[0.2'', 0.8'']$.

\paragraph{Sampler settings.} Every search uses \texttt{Nautilus} on one GPU with
\texttt{n\_batch} $= 10$ (the command-line default, applied to every stage; \texttt{--n\_batch}
overrides it) and one core. All other \texttt{Nautilus} settings are the library defaults
(\texttt{enlarge\_per\_dim} $= 1.1$, \texttt{split\_threshold} $= 100$, \texttt{n\_networks} $= 4$,
\texttt{f\_live} $= 0.01$, \texttt{n\_eff} $= 500$, \texttt{n\_shell} $= 1$, no fixed seed).
JAX runs in double precision. The number of live points grows by a fixed increment for each band
beyond the first, scaled to the parameters each band adds to that stage. Early stages are capped with
\texttt{n\_like\_max}; the no-subhalo baseline and the subhalo searches are not, because a search
stopped before convergence returns a biased evidence.

Table~\ref{tab:pipeline_stages} summarizes the stages; $N$ is the number of bands.

\begin{table}
    \centering
    \footnotesize
    \setlength{\tabcolsep}{4pt}
    \caption{\label{tab:pipeline_stages}Stages of the real-data pipeline, for $N$ bands. Free
    parameters are given for one and for two bands. Every stage uses \texttt{n\_batch} $= 10$. The
    position-penalty threshold is either fixed or computed from \texttt{source\_lp[1]}.}
    \begin{tabular}{lcccc}
        \hline
        Stage & Free (1 / 2) & $n_\mathrm{live}$ & $n_\mathrm{like,max}$ & Pos.\ threshold \\
        \hline
        \texttt{source\_lp[1]}  & 17 / 30 & $200 + 80(N-1)$  & $2 \times 10^5$ & fixed, $0.3''$ \\
        \texttt{source\_pix[1]} & 10 / 16 & $150 + 50(N-1)$  & $4 \times 10^5$ & computed \\
        \texttt{source\_pix[2]} & 3 / 6   & $75 + 25(N-1)$   & $5 \times 10^4$ & none \\
        \texttt{light[1]}       & 6 / 12  & $300 + 100(N-1)$ & $2 \times 10^5$ & none \\
        \texttt{mass\_total[1]} & 8 / 11  & $200 + 50(N-1)$  & none            & computed \\
        \texttt{subhalo[1]}, per cell & 11 / 14 & 600      & none            & computed \\
        \texttt{subhalo[2]}     & 11 / 14 & 600              & none            & computed \\
        \hline
    \end{tabular}
\end{table}

\subsection{\label{sc:pipeline_source_lp}\texttt{source\_lp[1]}: parametric source}

The lens light, lens mass and a parametric source are fitted together.
\begin{itemize}
    \item \textbf{Lens mass (shared):} an elliptical \texttt{Isothermal}. Centre: Gaussian priors,
    mean $0''$, $\sigma = 0.1''$. Elliptical components: truncated Gaussian priors, mean 0,
    $\sigma = 0.3$, truncated to $[-1, 1]$. Einstein radius: uniform on $[0'', 3.5'']$ (zero to the
    mask radius). The mass centre is free and is not tied to any band's lens-light centre.
    \item \textbf{External shear (shared):} $\gamma_1$ and $\gamma_2$ uniform on $[-0.3, 0.3]$.
    \item \textbf{Lens light (per band):} an MGE of two bases of 30 linear Gaussians each. The
    Gaussian widths are fixed and log-spaced from one tenth of the band's pixel scale to the mask
    radius. Both bases share one centre, with uniform priors of $\pm 0.1''$ about the band's mask
    centre. Each basis has its own elliptical components, with uniform priors on $[-0.5, 0.5]$. This
    gives six non-linear parameters per band.
    \item \textbf{Source (per band):} an MGE of one basis of 30 linear Gaussians, with widths
    log-spaced from one tenth of the pixel scale to $1.0''$. Centre: Gaussian priors on the mask
    centre with $\sigma = 0.6''$. Elliptical components: truncated Gaussian priors as for the mass.
    This gives four non-linear parameters per band.
    \item \textbf{Alignment:} free for every non-reference band, with the priors of
    Section~\ref{sc:pipeline_general}.
    \item \textbf{Position penalty:} a fixed source-plane threshold of $0.3''$.
\end{itemize}
A separate lens-light stage is not run first, because the lens is a single isolated galaxy.

\subsection{\label{sc:pipeline_source_pix_1}\texttt{source\_pix[1]}: first pixelized source}

Each band's parametric source is replaced by a pixelized source.
\begin{itemize}
    \item \textbf{Source (per band):} a Delaunay mesh whose image-plane vertices come from a Hilbert
    image mesh. The Hilbert mesh has 1500 points at a $0.03''$ pixel scale and 1250 at $0.06''$, with
    a weight power of 3.5 and a weight floor of 0.01. Its adapt data is the band's
    lens-subtracted source signal-to-noise map from \texttt{source\_lp[1]}, uncapped. Thirty further
    points are placed on a circle just outside the mask (radius: mask radius plus half a pixel), and
    the 30 Delaunay pixels they form are zeroed so that they carry no source flux.
    \item \textbf{Regularization (per band):} \texttt{AdaptSplit}, driven by the same signal-to-noise
    map. Inner and outer coefficients: log-uniform priors on $[10^{-6}, 10^6]$. Signal scale: uniform
    on $[0, 1]$. This gives three parameters per band.
    \item \textbf{Over-sampling of the pixelized source (per band):} sub-grid size 4 in pixels whose
    source signal-to-noise exceeds 3, and 2 elsewhere.
    \item \textbf{Lens mass (shared):} \texttt{Isothermal}, free with the \texttt{source\_lp[1]}
    priors.
    \item \textbf{External shear (shared):} free with the \texttt{source\_lp[1]} priors.
    \item \textbf{Lens light (per band):} fixed at its \texttt{source\_lp[1]} maximum-likelihood
    geometry.
    \item \textbf{Alignment:} free with the \texttt{source\_lp[1]} priors.
    \item \textbf{Position penalty:} threshold computed from \texttt{source\_lp[1]}.
\end{itemize}

\subsection{\label{sc:pipeline_source_pix_2}\texttt{source\_pix[2]}: refined pixelized source}

Each band's image mesh, adapt images and pixelization over-sampling are rebuilt from its
\texttt{source\_pix[1]} reconstruction. The adapt data fed to the image mesh is capped at a
signal-to-noise of 3, so that the brightest source pixels do not take most of the mesh. The adapt
images used by the regularization and the over-sampling remain uncapped. The datasets and adapt
images built here are reused by every later stage.
\begin{itemize}
    \item \textbf{Free:} only each band's \texttt{AdaptSplit} parameters, with the priors above.
    \item \textbf{Fixed:} the lens mass, shear and alignment at their \texttt{source\_pix[1]}
    maximum-likelihood values, and the lens light at its \texttt{source\_lp[1]} values.
    \item \textbf{Position penalty:} none.
\end{itemize}
The bands share no free parameter, so this joint fit is equivalent to independent per-band fits.

\subsection{\label{sc:pipeline_light}\texttt{light[1]}: lens light}

The lens light is refitted against the pixelized sources.
\begin{itemize}
    \item \textbf{Lens light (per band):} a fresh MGE of two bases of 30 linear Gaussians, with the
    same widths as in \texttt{source\_lp[1]}. Centre: uniform priors of $\pm 0.1''$ about that band's
    \texttt{source\_lp[1]} lens-light centre. Elliptical components: truncated Gaussian priors, mean
    0, $\sigma = 0.3$, truncated to $[-1, 1]$; these differ from the uniform $\pm 0.5$ priors of
    \texttt{source\_lp[1]}.
    \item \textbf{Fixed:} each band's pixelized source (mesh and regularization) from
    \texttt{source\_pix[2]}, and the lens mass, shear and alignment from \texttt{source\_pix[1]}.
    \item \textbf{Position penalty:} none.
\end{itemize}
The bands again share no free parameter.

\subsection{\label{sc:pipeline_mass}\texttt{mass\_total[1]}: total mass and the no-subhalo baseline}

The shared lens mass becomes an elliptical \texttt{PowerLaw}.
\begin{itemize}
    \item \textbf{Lens mass (shared):} \texttt{PowerLaw}. Its centre, elliptical components and
    Einstein radius take the \texttt{source\_pix[1]} \texttt{Isothermal} priors. Density slope:
    uniform on $[1.5, 3.0]$.
    \item \textbf{External shear (shared):} free with its original priors.
    \item \textbf{Alignment:} free with its original priors.
    \item \textbf{Lens light (per band):} fixed from \texttt{light[1]}.
    \item \textbf{Source (per band):} fixed from \texttt{source\_pix[2]} (mesh, image-plane mesh
    and regularization).
    \item \textbf{Position penalty:} threshold computed from \texttt{source\_lp[1]}.
    \item \textbf{Sampling:} uncapped, because the evidence of this fit is the no-subhalo baseline
    of the detection.
\end{itemize}
This fit is not repeated before the subhalo search. A refit without a subhalo, on the same analyses
and with the same priors, would reproduce it, because chaining reuses the original priors.

\subsection{\label{sc:pipeline_subhalo}Subhalo search}

The two subhalo stages compose their models from \texttt{mass\_total[1]}. The lens mass, shear and
alignment are free with the same priors as in \texttt{mass\_total[1]}. The lens light and pixelized
source of each band are fixed at their \texttt{mass\_total[1]} values. The fits run on the very
analysis objects used by \texttt{mass\_total[1]}: the same datasets, adapt images and position
penalty. They therefore differ from the baseline only by one subhalo, shared by all bands.

\paragraph{Subhalo model.} A spherical NFW halo at the lens redshift, whose concentration follows
the mass--concentration relation of \citet{Ludlow2016} (\texttt{NFWMCRLudlowSph}). Its only mass
parameter is $M_{200}$, with a log-uniform prior on $[10^6, 10^{11}]\,\mathrm{M}_\odot$. The lens
and source redshifts in the profile are fixed to those of the system.

\paragraph{\texttt{subhalo[1]}: grid search.} The subhalo centre is searched on a $2 \times 2$ grid
over $(-3'', 3'')$ in $y$ and $x$. Each cell is a separate joint fit of all bands, with uniform
priors on its $3'' \times 3''$ cell for the centre. \citet{Nightingale2024} used $5 \times 5$ over
the same region; $2 \times 2$ is chosen to cut the cost of the grid. The cells run one after another
on one GPU. The grid can be changed with \texttt{--grid\_steps} and \texttt{--grid\_dimension}.

\paragraph{\texttt{subhalo[2]}: refinement.} The subhalo is refitted with its centre free across the
image plane. The centre's priors are Gaussian, centred on the maximum-likelihood subhalo centre of the
highest-\emph{evidence} grid cell, with $\sigma$ equal to half a cell width ($1.5''$ for the
$2 \times 2$ grid). $M_{200}$ keeps its full log-uniform prior. The best cell is chosen by evidence
rather than likelihood, because the evidence is the detection statistic.

\paragraph{Detection statistic.} Two increases relative to \texttt{mass\_total[1]} are reported:
\begin{align}
    \Delta \ln Z_\mathrm{grid} &= \ln Z_\mathrm{best\ cell} - \ln Z_{\texttt{mass\_total[1]}}, \\
    \Delta \ln Z_\mathrm{refine} &= \ln Z_{\texttt{subhalo[2]}} - \ln Z_{\texttt{mass\_total[1]}}.
\end{align}
Both depend on the centre prior through its Occam factor. A $3'' \times 3''$ cell is $\ln 4 \approx
1.4$ denser at the subhalo than a uniform prior over the whole $6'' \times 6''$ region. Taking the
best of $N_\mathrm{cells}$ cells is a look-elsewhere selection, so
$\Delta \ln Z_\mathrm{grid} - \ln N_\mathrm{cells}$ approximates the evidence for a subhalo whose
position is unknown. The refined value is typically about 0.5 below the grid value, but its prior is
centred using the same data, which favours the subhalo model. The two are read together. Absolute
evidences of different band sets are never compared, because they are evidences of different data.

\subsection{\label{sc:pipeline_mock}Differences for mock data}

The mock data are fitted with the same pipeline up to and including \texttt{mass\_total[1]}. Every
stage up to that point has the same model, priors and sampler settings, so the two pipelines share
those fits. The differences are:
\begin{itemize}
    \item \textbf{No grid search.} \texttt{subhalo[1]} is not run, because on a simulation the
    subhalo's position is known.
    \item \textbf{Truth-centred subhalo prior.} In \texttt{subhalo[2]} the subhalo-centre priors are
    Gaussians centred on the true subhalo centre of the simulation ($(-0.10'', -1.90'')$, read from
    \texttt{info.json}), with $\sigma = 1.5''$ (\texttt{--centre\_sigma}). This is the width of the
    real-data \texttt{subhalo[2]} prior with the $2 \times 2$ grid. $M_{200}$, the sampler settings
    ($n_\mathrm{live} = 600$, uncapped, \texttt{n\_batch} $= 10$) and every other prior are unchanged.
    \item \textbf{One detection statistic},
    $\Delta \ln Z = \ln Z_{\texttt{subhalo[2]}} - \ln Z_{\texttt{mass\_total[1]}}$. Narrowing the centre
    prior by a factor two raises it by about $2 \ln 2 \approx 1.4$ while the prior is much wider
    than the posterior. This shift is the same in every band set, so the single- versus joint-band
    comparison does not depend on $\sigma$. The \texttt{no\_subhalo} control constrains no centre and
    barely responds to $\sigma$, so it is compared with \texttt{subhalo\_m10} at the same $\sigma$
    only. The values are not directly comparable with the real-data grid values.
    \item \textbf{Positions.} The multiple-image positions are solved from the true lens model of the
    simulation, so for \texttt{subhalo\_m10} they include the subhalo's small image-position shift.
    \item \textbf{Data inspection.} The simulated data have no contaminants, and every band uses the
    simulation's $3.5''$ mask about $(0'', 0'')$. The inspection step for real data is not needed.
\end{itemize}
